\section{Related work, implications, and limitations}\label{sec:related}
\paragraph{Information and investigative decisions.}
Blackwell's comparison of experiments \cite{blackwell1953} supplies a classical framework for ordering information by decision performance. CRTT fixes the decision contract and asks whether a common admissible action exists in each observed fiber; it does not claim a new general ordering of statistical experiments. Its exact finite observation-tree recurrence is an application of dynamic programming \cite{bellman1957}. The complexity of optimal decision-tree construction is a separate issue \cite{hyafil1976}; our exact synthesis theorem asserts no polynomial-time bound. Query learning \cite{angluin1988} and adaptive submodularity \cite{golovin2011} study acquisition strategies under different assumptions and objectives. The XOR witness in Section~3 makes clear why unrestricted greedy acquisition is not justified here. The threshold result in Section~2 likewise belongs to a restricted rent-versus-discovery model, not a universal online optimizer.

\paragraph{Learning from failure and verification.}
The conversion of failure into reusable constraints has established precedents: Soar learns procedural knowledge from impasse resolution \cite{soar1987}, and GRASP learns constraints from conflicts in propositional search \cite{grasp1999}. Certified SAT proof checking separates production from validation \cite{lrat2017}. CRTT specializes obstruction extraction to failed morphism rows, proves the resulting rule's scope and coverage preservation, and accounts for declared lookup, probing, and compilation charges. Soundness of a learned rule alone does not entail a profitable policy. Proof-carrying code \cite{necula1997} provides the general architectural precedent for checking a supplied proof against a fixed policy; in CRTT the corresponding obligation includes the denotation of unfinished search, comparative cost evidence and a completion reserve.

\paragraph{Reflection and changing control.}
Procedural reflection \cite{smith1984} makes computational mechanisms accessible to further computation. Rational metareasoning \cite{russell1991} evaluates computational actions by their contribution to decisions and their cost. Dynamic software updating \cite{stoyle2007} treats the correctness of changes to running programs. G\"odel machines \cite{schmidhuber2007} and formal models of rational self-modification \cite{everitt2016} address broader classes of self-change. CRTT does not subsume these frameworks: the checker, provider, contract and base tariffs are fixed, and the changeable policies have an explicitly finite syntax. Its concrete contribution is the composition of constructive transfer discovery, paid obstruction learning, continuation-preserving replacement and funded closure, rather than priority for reflection, conflict learning or proof-carrying verification considered separately.

\paragraph{Companion studies.}
The author's earlier studies of observability and retained memory in Rule~110 \cite{strand110observer2026,strand110memory2026} and decision-relevant information in finite channels \cite{strandcontrol2026} motivate the distinction between available information and usable control. They are neither imported proof dependencies nor empirical validation of CRTT. The supplement preserves the complete companion bibliography.

\paragraph{What the certificate producer must supply.}
The input boundary is proof-carrying, not a universal certificate constructor. Some certificates are computed inside the stack: the finite sum of job budgets $P(C,f)$ gives a sufficient reserve and hence $0\le b(c)\le R$ with $R\ge P(C,f)$ (Section~6.1). For a supplied funded plan of Section~5.1, \rc{initialBounds} constructs $[k,k+\lambda+R]$; \rc{exactRequest} evaluates that plan's modeled bill at a prepaid charge of its reservation plus a positive surcharge, producing a point interval. An absolute cost certificate embeds into a singleton joint row via \rc{fromAbsolute}. These are constructive but need not be cheap or tight. The closed fixture in Section~8.2 uses a synthesized acquisition tree and exact bounds proved by finite evaluation. In general, however, a useful old-cost lower bound, tight new-cost upper bound, covering joint rows, true-hypothesis membership, and factories indexed by the actual state must be provided. Evaluating exact bills can itself require completion; admission avoids repeating that work but does not eliminate it. There is no general procedure here for discovering such economical evidence or bounding its real production time. Thus closure certifies a conditional composition, not an efficient universal discovery algorithm.

\paragraph{Scope and research implications.}
The guarantees are relative to finite carriers, deterministic observations and sound witnesses. The model neither establishes empirical truth nor automatically prices proof construction in physical resources. The release does not migrate arbitrary program state, rewrite its kernel, or establish convergence of unbounded schedules. Typed in-memory provenance is not crash-safe persistence. The finite witnesses establish no generalization to neural agents or new task families. No exhaustive novelty comparison or asymptotic speedup theorem is claimed.

The epistemic interpretation is that a result of investigation can become a verified operator of subsequent investigation, without conflating knowledge of a method's correctness with knowledge of its utility. Negative certificates, partial cost evidence and continuation equivalence make different aspects of that process explicit. CRTT demonstrates, in a specific executable model, how such operators and their control policies can be revised while the obligation to finish the same root investigation remains invariant. Extending that contract to empirical domains is a separate research problem, not a corollary of the release theorem.
